Wir rechnen mit der Druckzunahme $\Delta p = \poeO$ gegenüber dem unkomprimierten Zustand. Mit der Kompressibilität von Öl $\chi = \chiOO$ (FoTaBe S.~188) und $|\Delta V| = \chi\,V\,\Delta p$ (FoTaBe S.~164) ist
\begin{align}
 V &= \VoeF \\
   &= \frac{\dVo}{\chiO\cdot\poe} \\
   &= \Voe \approx \result{\VoelP{0}{2}}
\end{align}
Da sich das Volumen nur um etwa ein Promille ändert, spielt es keine Rolle, ob $V$ das komprimierte oder das unkomprimierte Volumen ist. Ist mit \poeO\ der absolute Druck gemeint, ist die Druckzunahme gegenüber dem Normdruck $p_\mathrm{n} = \pnoO$ etwas kleiner, und man erhält $V = \VoaF \approx \VoalP{0}{2}$.
